% 本文件由 LaTeX 排版 Agent 根据有效 translation.json 自动生成。
% author=Clement of Rome; work_id=1010; title=致格林多人书（克肋孟）

\chapter{致格林多人书（克肋孟）}

% structure: p0001 source=h1 target=omitted reason=page-title-already-rendered-by-parent

% structure: p0002 source=h2 target=section reason=relative-heading-rank; base=h2

\section{第一章 问候。在分裂爆发前对格林多人的赞美。}

旅居罗马的天主的教会，致旅居于格林多的天主的教会，就是那些因天主的旨意，藉我们的主耶稣基督蒙召而成圣的人：愿恩宠与平安，由全能的天主藉耶稣基督，多多地赐予你们。\par

亲爱的弟兄们，由于我们自己遭遇的突发且接连不断的灾祸，我们感到在关注你们所咨询的问题上有些迟延；尤其是那件可耻可恶的叛乱，完全为天主的选民所憎恶，是由几个轻率自信的人煽动到如此疯狂的地步，致使你们那值得普世爱戴的尊贵而显赫的名声遭受严重损害。因为有谁曾在你们中间短暂居住，而没有发现你们的信德既坚固又富有德行？有谁不钦佩你们在基督内虔诚的节制与适中？有谁不宣扬你们惯常的慷慨好客？又有谁不为你们完善而稳固的知识而欣喜？因为你们行事从不偏私，遵行天主的诫命，服从那些治理你们的人，对你们中间的长老给予一切合宜的尊敬。你们嘱咐青年人要持守清醒庄重的心，教导你们的妻子以无可指摘、合宜而纯洁的良心行一切事，按本分爱自己的丈夫；你们也教导她们在服从的规则中生活，合宜地管理家务，在各方面都显明有节制。\par

% structure: p0005 source=h2 target=section reason=relative-heading-rank; base=h2

\section{第二章 对格林多人的赞美（续）。}

此外，你们全都以谦卑为荣，毫不自高自大，宁愿服从而不强求服从，更愿施予而不愿领受。\BibleRef{宗 20:35} 你们满足于天主为你们预备的一切，留心听祂的话语，内心充满祂的道理，祂的受苦摆在你们眼前。如此，深厚的平安丰丰富富地赐予你们众人，你们行善的渴望永无止境，圣神的充沛浇灌也临于你们众人。你们充满神圣的意念，以真诚的热忱和虔诚的信心，向全能的天主伸出双手，恳求祂怜悯你们，若你们犯了任何不自觉的过失。你们昼夜挂念整个弟兄团体，\BibleRef{伯前 2:17} 为使天主的选民能以怜悯和良心得救。你们真诚无伪，彼此不记怨。任何派别和分裂在你们眼中都是可憎的。你们为邻人的过犯哀恸，将他们的缺欠视为自己的。你们从不吝啬任何善行，\BibleQuote{准备行各种善事}\BibleRef{铎 3:1}。你们因全然有德而虔诚的生活而受装饰，在敬畏天主中行一切事。主的诫命和典章写在你们心版上。\BibleRef{箴 7:3}\par

% structure: p0007 source=h2 target=section reason=relative-heading-rank; base=h2

\section{第三章 因嫉妒和竞争引发分裂后，格林多教会的不幸状况。}

各样尊荣和幸福曾赐予你们，然后经上所记载的应验了：\BibleQuote{我所爱的人吃了喝了，长大肥胖，便踢跳。}\BibleRef{申 32:15} 由此涌出竞争与嫉妒、纷争与叛乱、迫害与混乱、战争与囚禁。于是卑贱者起来反对尊贵者，无名之辈反对有名之士，愚昧者反对智慧者，年轻人反对年长者。为此，公义与平安如今远离你们，因为人人都放弃了敬畏天主，在信仰上眼目昏花，既不遵行祂所命的典章，也不按基督徒当行的样式生活，却随从自己邪恶的情欲，重新拾起不义和邪恶的嫉妒，死亡就是藉此进入世界的。\BibleRef{智 2:24}\par

% structure: p0009 source=h2 target=section reason=relative-heading-rank; base=h2

\section{第四章 自古已从这根源流出许多邪恶。}

因为经上记载：\BibleQuote{过了一些日子，加音带了田地的出产献祭给天主；亚伯尔也带了羊群中头生的和脂肪来献。天主看顾了亚伯尔和他的祭品，却没有看顾加音和他的祭品。加音非常忧愁，脸色也沉了下来。天主对加音说：你为什么忧愁，脸色为什么沉下来？你若行得好，岂不蒙悦纳？若行得不好，罪就伏在门前；它必恋慕你，你却要制伏它。加音对他的弟弟亚伯尔说：我们到田间去吧。当他们还在田间的时候，加音起来攻击他的弟弟亚伯尔，把他杀了。}\BibleRef{创 4:3-8} 弟兄们，你们看，嫉妒和嫉恨如何导致谋杀兄弟。也因嫉妒，我们的祖先雅各伯逃避他兄弟厄撒乌的面。嫉妒使若瑟受迫害至死，并成为奴仆。\EditorialNote{创世纪第三十七章} 嫉妒迫使梅瑟逃避埃及王法郎的面，当他听到同胞说：\BibleQuote{谁立了你作我们的首领和审判官？难道你要杀我，像杀那埃及人一样吗？}\BibleRef{出 2:14} 因嫉妒，亚郎和米黎盎不得不住在营外。\BibleRef{户 12:14} 嫉妒使达堂和阿彼兰因他们鼓动叛乱反对天主的仆人梅瑟而活活堕入阴府。\BibleRef{户 16:33} 因嫉妒，达味不仅遭受外邦人的仇恨，也被以色列王撒乌耳迫害。\par

% structure: p0011 source=h2 target=section reason=relative-heading-rank; base=h2

\section{第五章 同样出自这根源的邪恶在近期也发生不少。伯多禄和保禄的殉道。}

但不必详论古代的榜样，让我们来到最近的灵性英雄中。让我们选取我们这一代所提供的高贵典范。由于嫉妒和争竞，［教会］中最伟大、最正义的支柱曾受迫害而被处死。让我们将卓越的宗徒摆在我们眼前。\Person{伯多禄}，由于不义的嫉妒，忍受了不止一次两次，而是许多次的劳苦；当他最终殉道后，离世前往他应得的荣耀之处。由于嫉妒，\Person{保禄}也获得了坚忍的赏报，他曾七次被囚、被迫逃亡、被石头砸。他在东方和西方宣讲后，因他的信德获得了卓越的名声，向全世界教导了正义，来到了极西之地，并在总督手下殉道。这样，他离开了世界，进入了圣所，成了忍耐的杰出榜样。\par

% structure: p0013 source=h2 target=section reason=relative-heading-rank; base=h2

\section{第六章 继续。其他几位殉道者。}

在这些度圣洁生活的人身上，还要加上一大群蒙选者，他们因嫉妒忍受了许多凌辱和折磨，为我们提供了最卓越的榜样。因着嫉妒，那些妇女——\Person{达纳伊德}和\Person{迪尔刻}——受迫害，在遭受可怕而难以言喻的折磨后，以坚定完成了她们的信德之赛程，虽身体软弱，却获得了高贵的赏报。嫉妒使妻子与丈夫疏离，改变了我们始祖\Person{亚当}的话：\BibleQuote{这才是我的骨中之骨，肉中之肉。} \BibleRef{创 2:23} 嫉妒和纷争推翻了伟大的城市，铲除了强大的民族。\par

% structure: p0015 source=h2 target=section reason=relative-heading-rank; base=h2

\section{第七章 劝告悔改。}

亲爱的，我们写这些给你们，不仅是为劝勉你们尽本分，也是为提醒我们自己。因为我们是在同一个竞技场上奋斗，同样的争战也分配给我们二人。因此，让我们放弃虚浮无益的挂虑，接近我们圣召那光荣可敬的规范。让我们注意在造我们的主眼中看为美好、悦纳和可喜的事。让我们坚定地注视基督的血，看那血在天主眼中是何等宝贵：它为我们得救而倾流，将悔改的恩宠摆在整个世界面前。让我们回顾已过的每一个时代，并得知从一代到一代，主已为所有愿意归向祂的人预留了悔改的地方。\Person{诺厄}宣讲了悔改，凡听从他的人都得了救。\Person{约纳}宣告了\Place{尼尼微}人的毁灭；\EditorialNote{约纳书第三章}但他们悔改了自己的罪，藉祈祷平息了天主的义怒，获得了救恩，虽然他们是［对天主的盟约而言的］外邦人。\par

% structure: p0017 source=h2 target=section reason=relative-heading-rank; base=h2

\section{第八章 关于悔改的继续。}

天主恩宠的仆役曾藉圣神论及悔改；万有的主自己也以誓言论及悔改说：\BibleQuote{我指着我的永生说，主说：我不愿意恶人死亡，却愿意他悔改；} \BibleRef{则 33:11} 此外又加上了这恩慈的宣告：\BibleQuote{以色列家啊，你们要悔改你们的罪过。} \BibleRef{则 18:30} 对我子民的儿女说：即使你们的罪从地上高及天上，即使它们比朱红更红\BibleRef{依 1:18}，比麻衣更黑，但若你们全心归向我，说：父啊！我必应允你们，如同应允圣洁的民。在另一处祂这样说：\BibleQuote{你们要洗濯，自洁；从我眼前除掉你们灵魂的邪恶；停止作恶，学习行善；寻求公正，解救受压迫者，为孤儿伸冤，为寡妇主持公道；然后我们来辩论。祂说：你们的罪虽像朱红，我必使他们白如雪；虽红如丹颜，我必使他们白如羊毛。你们若乐意听从我，必吃地上的美物；若不听从，悖逆不听话，刀剑必吞灭你们，因为这是主亲口说的。} \BibleRef{依 1:16} 因此，愿祂所有蒙爱的人都有分于悔改，祂以其大能的旨意确立了［这些宣告］。\par

% structure: p0019 source=h2 target=section reason=relative-heading-rank; base=h2

\section{第九章 圣人们的榜样。}

因此，让我们服从祂卓越而光荣的旨意；恳求祂的怜悯和慈爱，同时舍弃一切徒劳的努力、纷争和引人死亡的嫉妒，让我们转向并投靠祂的慈悲。让我们坚定地默观那些完美地事奉祂卓越光荣的人。让我们以\Person{哈诺客}为例：他因顺服而被称为义，被提升天，死亡从未临到他。\Person{诺厄}因忠信而被寻见，藉他的职务向世界宣讲了新生；主拯救了那些齐声进入方舟的动物。\par

% structure: p0021 source=h2 target=section reason=relative-heading-rank; base=h2

\section{第十章 上述内容的继续。}

亚巴郎，被称为\Quote{朋友}，因他服从了天主的言语，而被视为有信德的人。他因着服从，离开了自己的故乡、本族和父家，为要舍弃一片小地、一个弱小的家族和微不足道的家，从而继承天主的恩许。因为天主曾对他说：\Quote{你离开你的故乡、你的本族和你的父家，往我指给你的地方去。我必使你成为一个大民族，我必降福你，使你的名受显扬，你必成为福源。降福你的，我必降福他；咒骂你的，我必咒骂他；地上万民都要因你获得祝福。}\BibleRef{创 12:1} 又当他离开罗特时，天主对他说：\Quote{你举起眼来，由你所住的地方往东西南北观看；凡你所看见的地方，我都要赐给你和你的后裔，永远作为产业。我要使你的后裔多如地上的尘土，如果人能数清地上的尘土，你的后裔也能数清。}\BibleRef{创 13:14} 又［圣经］说：\Quote{天主领亚巴郎到外面，对他说：请你仰观苍天，数点星辰，你能够数清吗？你的后裔也将这样。亚巴郎相信了天主，天主就以此算为他的正义。}由于他的信德和好客，他在年老时得了一个儿子；又因着服从，他将儿子献为祭献，放在天主指示给他的一座山上。\par

% structure: p0023 source=h2 target=section reason=relative-heading-rank; base=h2

\section{第十一章　续·罗特}

由于罗特的好客与虔敬，当整个地区因硫磺与火而受罚时，他得以从索多玛得救，主借此表明：祂不遗弃那些寄望于祂的人，却将离弃祂的人交给惩罚与折磨。\BibleRef{伯后 2:6} 因为罗特的妻子与他一同出来，却心志不同，不肯与他同心（关于所赐的命令），就成了一个鉴戒，直到今日仍是一根盐柱。这事为叫众人知道，那些三心二意、不信天主大能的人，自招审判，并成为后世各代的征兆。\par

% structure: p0025 source=h2 target=section reason=relative-heading-rank; base=h2

\section{第十二章　信德与好客的赏报·辣哈布}

由于辣哈布妓女的信德和好客，她得到了救恩。因为当农的儿子若苏厄派遣探子往耶里哥去时，那地的国王知道他们来窥探地方，就派人去捉拿他们，要将他们处死。但好客的辣哈布接待了他们，把他们藏在房顶的麻秸里。当国王所派的人来到说：\Quote{有几个人到你这里来，要窥探这地方；把他们交出来，因为王如此命令。}她回答他们说：\Quote{你们所找的那两个人，曾到我这里来，但很快就离开了，已经走了。}这样她没有向那些人暴露探子。然后她对那些人说：\Quote{我确实知道，主你们的天主已将这城赐给你们，因为你们害怕和恐惧已临到居民身上。所以当你们攻取这城时，请保全我和我父家的性命。}他们对她说：\Quote{必照你对我们说的话办理。因此，你一知道我们临近，就要将你全家聚集在你的屋顶下，他们必得保全；凡在外面的人，必灭亡。}此外，他们给了她一个记号：她要从家中挂出一根朱红线绳。他们以此表明，救赎应藉着主的血，流向所有相信并寄望于天主的人。亲爱的诸位，你们看，在这女人身上不仅有信德，而且有预言之恩。\par

% structure: p0027 source=h2 target=section reason=relative-heading-rank; base=h2

\section{第十三章　劝谕谦卑}

因此，弟兄们，让我们存心谦卑，除去一切傲慢、骄傲、愚妄和愤怒的情绪；并依照经上记载而行（因为圣神说：\Quote{智者不要夸耀自己的智慧，勇士不要夸耀自己的勇力，富人不要夸耀自己的财富；凡要夸耀的，应因主而夸耀，专心寻求祂，并施行审判和正义。}），尤其要铭记主耶稣的话，祂教导我们温良和忍耐。因为祂这样说：\Quote{你们要慈悲，好蒙慈悲；宽恕，好被宽恕；你们怎样待人，人也怎样待你们；你们怎样判断，人也怎样判断你们；你们怎样仁慈，人也怎样仁慈待你们；你们用什么尺度量人，人也用什么尺度量你们。}让我们以这条诫命和这些规则建立自己，好使我们满怀谦卑，服从祂的圣言而行。因为圣言说：\Quote{我垂顾谁呢？惟有垂顾那温顺、平安、敬畏我言语的人。}\BibleRef{依 66:2}\par

% structure: p0029 source=h2 target=section reason=relative-heading-rank; base=h2

\section{第十四章　我们该服从天主，而非制造纷争的人}

因此，弟兄们，正当而神圣的事是：服从天主，而非追随那些因骄傲和纷争而成为可憎嫉妒之头目的人。因为如果我们轻率地顺从那些煽动争端和骚乱、企图使我们远离美善之人的意念，必将招致不小的伤害，反而是极大的危险。让我们以造物主的仁慈和良善为榜样，彼此以善意相待。因为经上记载：\Quote{心地善良的必安居乐业，无罪的人要存留地上，但恶人要从地上消灭。}\BibleRef{箴 2:21} 又［圣经］说：\Quote{我看见恶人高升，耸立如黎巴嫩的香柏；我再经过，已不见他；我寻找他的地方，也寻不着。你应持守纯善，注视正直，因为和平的人必留下后裔。}\par

% structure: p0031 source=h2 target=section reason=relative-heading-rank; base=h2

\section{第十五章　我们当依附培养和平的人，而非那些只是假装如此的人}

让我们因此依附那些以虔诚培育和平的人，而非那些虚伪地声称渴望和平的人。因为\WorkTitle{圣经}在某处说：\BibleQuote{这百姓用嘴唇尊敬我，心却远离我。}又说：\BibleQuote{他们用口祝福，心却咒诅。}又说：\BibleQuote{他们用口爱祂，用舌头向祂说谎；他们的心不向祂正直，在祂的盟约上也不忠诚。}愿说谎的嘴唇沉默，[并且]\BibleQuote{愿上主消灭一切说谎的嘴唇和夸大之舌，就是那些说\InnerQuote{我们要夸大我们的舌头；我们的嘴唇是自己的，谁能作我们的主？}的人。因贫寒人的受压迫，因穷困人的叹息，我现在要起来，上主说：\InnerQuote{我要把他安置在稳妥之地，我要坦然对待他。}}\par

% structure: p0033 source=h2 target=section reason=relative-heading-rank; base=h2

\section{第16章 基督——谦卑的榜样。}

因为基督是属于那些谦卑的人，而非那些高举自己、践踏祂羊群的人。我们的主耶稣基督，天主威严的权杖，并未以骄矜傲慢的姿态降临，尽管祂本可如此，却以卑微的身份而来，正如圣神关于祂所预言的。因为祂说：\BibleQuote{上主，谁相信了我们的报告，上主的臂膀又向谁显示了呢？我们在祂面前宣告了[我们的信息]：祂好似一个孩童，又像干涸之地的一根根苗；祂没有俊美，也没有光荣；我们看见祂的时候，祂没有仪容，没有美貌；但祂的形貌没有出众之处，甚至比一般人的形貌更为低下。祂是一个受鞭打、受苦楚的人，熟悉忍受痛苦；因为祂的面容被转开，祂被轻视，不被尊重。祂承担了我们的罪过，为我们而忧伤；我们却以为祂是因自己的缘故而受劳苦、受鞭打、受磨难。然而祂是因我们的过犯而被刺伤，因我们的罪孽而被压伤。祂所受的惩罚使我们获得平安，因祂的创伤我们得了痊愈。我们都像羊一样迷了路，各人行走自己的道路；上主却把祂交出来，为我们的罪而受苦，祂在受苦中却不开口。祂被牵去宰杀，如同羊羔，又像母羊在剪毛的人前缄默无声，祂也不开口。在祂的屈辱中，祂的审判被剥夺了；谁能述说祂的世代？因为祂的生命从地上被夺去。因我百姓的过犯，祂被带到死地。我必将恶人当作祂的坟墓，将富人当作祂的死亡，因为祂没有行过不义，口中也没有诡诈。上主喜悦用鞭打来洁净祂。你若为罪献上祭物，你的灵魂必要看见长久的后裔。上主喜悦解救祂灵魂的痛苦，给祂显示光明，以智慧塑造祂，使那一位正义者成义，祂善加服事众人；祂自己要承担他们的罪过。因此，祂必继承许多，并分取强者的战利品，因为祂的灵魂被交付于死亡，祂被列于罪犯之中，祂承担了许多人的罪，并为他们的罪而被交出。}又说：\BibleQuote{我是虫，而不是人，是人的耻辱，是百姓所轻视的。凡看见我的都嘲笑我，他们嘴唇说话，摇头说：他信赖上主，让上主解救他吧，让上主拯救他吧，既然上主喜爱他。}你们看，亲爱的，这是赐给我们的榜样；如果上主这样谦卑自己，那么，我们这些靠着祂而负起祂恩宠之轭的人，又该怎样做呢？\par

% structure: p0035 source=h2 target=section reason=relative-heading-rank; base=h2

\section{第17章 圣人——谦卑的榜样。}

让我们也效法那些披着山羊皮和绵羊皮\BibleRef{希 11:37}四处宣讲基督来临的人，我指的是先知中的厄里亚、厄里叟、厄则克耳，以及\WorkTitle{圣经}中同样受称赞的其他人。\Person{亚巴郎}特别受尊敬，被称为天主的朋友；然而他真诚地看重天主的光荣，谦卑地宣称：\BibleQuote{我不过是灰尘和灰烬。}\BibleRef{创 18:27}此外，关于约伯也这样记载：\BibleQuote{约伯是个正义、正直、敬畏天主、远离一切恶事的人。}\BibleRef{约 1:1}但他却控告自己说：\BibleQuote{没有人是不洁的，即使他的生命只有一天。}\BibleRef{约 14:4}\Person{梅瑟}在天主的全家中被称为忠信；藉着他，天主用瘟疫和磨难惩罚了埃及人。然而，尽管他如此受尊崇，他却不说高傲的话，当天主从荆棘中向他发出神谕时，他说：\BibleQuote{我是谁，竟使你派遣我？我是拙口笨舌的人。}又说：\BibleQuote{我不过像锅里的烟雾。}\par

% structure: p0037 source=h2 target=section reason=relative-heading-rank; base=h2

\section{第18章 达味——谦卑的榜样。}

但关于\Person{达味}，我们该说什么呢？他得了这样的见证，天主曾说：\BibleQuote{我找着一个合我心意的人，\Person{达味}，\Person{叶瑟}的儿子；我要用永远的仁慈傅油于他？} 然而这人对天主说：\BibleQuote{天主，求祢按祢的仁慈怜悯我，按祢丰厚的怜悯涂去我的过犯。求祢将我的过犯洗尽，洗净我的罪，因为我承认我的过犯，我的罪常在我面前。我唯独得罪了祢，在祢眼前行了恶，为叫祢在祢的言语上显为公义，在祢受审判时得胜。看，我是在过犯中怀的，在罪中我母亲怀了我。看，祢喜爱真理，将智慧和隐藏的事启示了我。祢用牛膝草洒我，我就洁净；祢洗涤我，我就比雪更白。祢使我听到喜乐和欢畅，使我卑微的骨头欢跃。求祢转面不看我的罪，涂去我的一切愆尤。天主，求祢在我内造一颗清洁的心，使我重有正直的灵。不要将我从祢面前抛弃，不要从祢那里收回祢的圣神。求祢还我祢救恩的喜乐，以祢统御的灵坚固我。我要将祢的道教导过犯的人，罪人必归向祢。天主，我救恩的天主，求祢救我脱离血债；我的舌头必要歌颂祢的公义。主，求祢开启我的唇，我的口要宣扬祢的赞美。因为祢若喜欢祭献，我必献上；但祢不喜爱燔祭。天主所要的祭献就是痛悔的精神，痛悔和谦卑的心，天主必不轻看。}\par

% structure: p0039 source=h2 target=section reason=relative-heading-rank; base=h2

\section{第十九章 效法这些榜样，让我们寻求和平。}

这样，伟大和杰出之人的谦卑与虔敬的顺服，不仅使我们，也使我们在前的所有世代变得更好，就是所有用畏惧和真理接受祂圣言的人。因此，既然我们有这么多伟大光荣的榜样摆在面前，让我们再次转向那从起初就放在我们面前的和平的实践；让我们坚定地仰望万有的圣父和创造者，紧紧依附祂那大能而极其伟大的和平恩赐和恩惠。让我们以理智默观祂，以我们灵魂的眼睛注视祂的容忍之旨意。让我们反思祂对祂所有受造物是多么远离愤怒。\par

% structure: p0041 source=h2 target=section reason=relative-heading-rank; base=h2

\section{第二十章 宇宙的和平与和谐}

诸天在祂的治理下运行，和平地服从祂。昼夜按祂所定的路线运行，互不阻碍。日、月并众星，遵照祂的命令，在他们限定的轨道内和谐运转，毫不偏离。大地按照祂的旨意，按时节为人类、牲畜和其上的一切生物丰盛地产出食物，从不迟延，也不改变祂所定的任何规律。深渊不可测之处，和下界的不可描述的安排，都被同样的规律约束。浩瀚无际的大海，因祂的作为汇聚成不同的盆地，从不越过它周围的界限，而是按祂所吩咐的行事。因为祂说：\BibleQuote{你只可到这里，不可越过；你狂傲的浪要到此为止。} \BibleRef{约 38:11} 那人不可渡过的海洋，以及它之外的世界，都由主同样的法令管理。春、夏、秋、冬的季节和平地相互交替。四方的风在适当的时候毫无阻碍地履行它们的服务。常涌的泉水，为享乐和健康而设，不断供应乳房来滋养人的生命。最微小的生物也在和平与协调中相遇。这一切伟大的创造者和万有的主安排了它们存在于和平与和谐中；而祂向众人施恩，但最丰盛地施予我们这些藉着我们的主\Person{耶稣基督}投奔祂怜悯的人，愿荣耀和威严归于祂，直到永远。阿们。\par

% structure: p0043 source=h2 target=section reason=relative-heading-rank; base=h2

\section{第二十一章 让我们服从天主，而非叛乱的制造者}

亲爱的，请注意，免得祂的许多恩慈导致我们所有人的定罪。[因为必须这样]除非我们行事为人与祂相称，并同心合意地做祂眼中美好和喜悦的事。因为[圣经]在某处说：\BibleQuote{上主的灵是烛火，洞察人的肺腑。} \BibleRef{箴 20:27} 让我们反思祂是多么临近，我们所有的思想和意念没有一样向祂隐藏。因此，我们不应该离开祂的旨意所分配给我们的岗位。我们宁愿得罪那些愚蠢、轻率、高傲、并以言语的骄傲自夸的人，而不得罪天主。让我们尊敬主\Person{耶稣基督}，祂的血为我们而倾流；让我们尊重那些管理我们的人；让我们尊敬我们中的长者；让我们在敬畏天主中培养年轻人；让我们引导我们的妻子走向善事。让她们表现出纯洁的可爱的习惯[在一切行为中]；让她们显露出温良的真诚气质；让她们通过说话的方式显明她们对舌头的掌控；让她们表达爱，不是偏爱这个胜过那个，而是对一切虔诚敬畏天主的人表现同等的爱。让你们的子女接受真正的基督徒训练；让他们学习谦卑在天主面前何等有益——纯净的爱情之灵能如何影响祂——祂的敬畏是何等卓越伟大，以及它如何拯救所有以纯洁的心行走在其中的人。因为祂是[人心]思念和欲望的探查者：祂的气息在我们内；当祂愿意时，祂会收回它。\par

% structure: p0045 source=h2 target=section reason=relative-heading-rank; base=h2

\section{第二十二章 这些劝诫由基督徒信仰所证实，该信仰宣告罪恶行为的悲惨}

现在那在基督内的信德确认了这一切劝诫。因为祂自己藉着圣神这样对我们说：\Quote{孩子们，来，听我言，我要教导你们敬畏上主。谁愿生活幸福，喜爱享见好日子？就要禁止舌头不出恶言，嘴唇不说欺诈的话。躲避罪恶，努力行善，寻求和平，追随不懈。上主的双目注视着义人，祂的耳朵[开启]垂听他们的祈祷。上主的面容对付作恶的人，要从地上除灭他们的纪念。义人呼求，上主听见了，就救他脱离一切患难。}\Quote{恶人的惩罚多，但仰望上主的人必被仁慈环绕。}\par

% structure: p0047 source=h2 target=section reason=relative-heading-rank; base=h2

\section{第二十三章 当谦卑，并相信基督将再来。}

极仁慈而恩惠的圣父对待敬畏祂的人怀有慈悲心肠，并以仁慈和爱心将祂的恩惠赐予那些以纯朴之心来到祂面前的人。为此，我们不要心怀二意；我们的灵魂也不可因祂那极其伟大而荣耀的恩赐而骄傲。但愿那所记载的话远离我们：\Quote{心怀二意、疑惑不定的人有祸了！他们说：\InnerQuote{这些事我们在我们祖先的时代就已听说；但是，看哪，我们已经老了，这些事一件也没有临到我们。}}你们愚昧的人哪！把自己与一棵树相比；例如葡萄树。首先，它落叶，然后发芽，接着长叶，然后开花；之后是青葡萄，随后是成熟的果实。你看，在很短的时间内，一棵树的果子就成熟了。的确，很快且突然地，祂的旨意必要成就，正如圣经也作证说：\Quote{祂必快来，决不迟延}；以及\Quote{上主必忽然进入祂的殿，就是你们所盼望的圣者。}\BibleRef{拉 3:1}\par

% structure: p0049 source=h2 target=section reason=relative-heading-rank; base=h2

\section{第二十四章 天主在大自然中不断向我们显示将有复活。}

我们要思想，亲爱的诸位，上主如何不断地向我们证明将来有复活，祂已藉着使主耶稣基督从死者中复活，使祂成为初果。我们要默想，亲爱的诸位，这不断发生的复活。昼夜向我们宣示复活。黑夜沉睡，白昼升起；白昼又离去，黑夜来临。让我们观看田间的果实，看谷粒是如何播种的。撒种者出去，把种子撒在地里；种子如此撒下，虽然落地时干燥赤裸，但渐渐消解。然后，藉着消解，天主眷顾的强有力力量又将它兴起，从一粒种子生出许多，并结出果实。\BibleRef{路 8:5}\par

% structure: p0051 source=h2 target=section reason=relative-heading-rank; base=h2

\section{第二十五章 凤凰是我们复活的象征。}

我们要思考那发生在东方之地，即在阿拉伯及其周围地区的奇妙迹象。有一种鸟叫作凤凰。这是独一无二的，活五百年。当它即将解体而死的时刻临近，它用乳香、没药和其他香料筑一个巢；时候一到，就进入巢中死去。但当肉体腐烂时，生出一种蠕虫，以死鸟的汁液为养分，长出羽毛。然后，当它强壮时，就抬起那个包含其父母骨头的巢，带着这些，从阿拉伯地进入埃及，到那称为赫利奥波利斯的城。并在光天化日之下，在众人眼前飞翔，将它们放在太阳的祭坛上；做完这事后，迅速返回原先的住处。祭司于是查阅日期记录，发现它正好在五百年结束时返回。\par

% structure: p0053 source=h2 target=section reason=relative-heading-rank; base=h2

\section{第二十六章 那么，我们将复活，正如圣经也作证。}

我们岂认为万物创造者将那些以良好信德的保证虔诚事奉祂的人再次兴起，是何等伟大奇妙的事？因为祂甚至藉着一只鸟向我们显示祂实现诺言的强大能力。因为[圣经]在某处说：\Quote{祢要使我起来，我必向祢承认}；又说：\Quote{我躺下，就睡了}；\Quote{我醒了，因为祢与我同在}；约伯又说：\Quote{祢要使我这承受了这一切的肉体兴起。}\BibleRef{约 19:25}\par

% structure: p0055 source=h2 target=section reason=relative-heading-rank; base=h2

\section{第二十七章 怀着复活的望德，让我们紧紧依附全能和全知的天主。}

我们既怀有这个望德，就让我们的灵魂紧紧依附那在祂的许诺上信实、在祂的审判上公义的那一位。祂命令我们不可说谎，祂自己更不会说谎；因为在天主，除了说谎，没有不可能的事。因此，让我们的信德再次被激发，并考虑万有都靠近祂。祂以祂权能的言语建立了万有，也以祂的言语可以推翻它们。\Quote{谁能对祂说：祢做了什么？或，谁能抵挡祂力量的大能？}祂愿意何时、如何，就何时、如何成就一切，祂所决定的事没有一件会消逝。\BibleRef{玛 24:35}万有在祂面前都是敞开的，没有事物能隐藏于祂的计谋。\Quote{诸天宣扬天主的荣耀，穹苍显示祂手的化工。日复一日发出言语，夜复夜传出知识。没有言语，没有话语，也听不到它们的声音。}\par

% structure: p0057 source=h2 target=section reason=relative-heading-rank; base=h2

\section{第二十八章 天主看见一切：因此让我们远离过犯。}

既然万物都被［天主］看见和听见，让我们敬畏祂，并弃绝那些出自邪恶欲望的恶行；如此，藉着祂的仁慈，我们就能免受将来的审判。因为，我们中谁能逃出祂大能的手呢？或者，哪个世界会接纳那些逃避祂的人呢？因为圣经在某处说：\BibleQuote{我往哪里去，在哪里能躲避祢的面容？我若升到天上，祢在那里；我若逃到地极，祢的右手也在那里；我若在深渊里铺床，祢的圣神也在那里。}那么，谁还能去哪里，或在哪里逃脱那包罗万有的祂呢？\par

% structure: p0059 source=h2 target=section reason=relative-heading-rank; base=h2

\section{第二十九章　让我们也以纯洁之心亲近天主。}

当至高者划分万民，当祂分散亚当的子孙时，祂按照天主众天使的数目，固定了万民的疆界。祂的子民雅各伯成了上主的产业，以色列成了祂嗣业的份。\BibleRef{申 32:8}在另一处，［圣经］说：\BibleQuote{看，上主从万民中为自己选取一个民族，如同人从禾场上取初果；从这个民族要产生至圣者。}\par

% structure: p0061 source=h2 target=section reason=relative-heading-rank; base=h2

\section{第三十章　让我们做那些中悦天主的事，并逃避祂所憎恶的事，好使我们蒙福。}

看见我们既是圣者的份子，就让我们做一切属于圣洁的事，避免一切诽谤、一切可憎和不洁的拥抱，连同一切酗酒、追求变化、一切可憎的情欲、可憎的奸淫和可憎的骄傲。因为［圣经］说：\BibleQuote{天主拒绝骄傲人，却赐恩宠给谦卑人。}那么，让我们紧贴那些由天主赐予恩宠的人。让我们以和谐与谦卑为衣，时常操练自制，远离一切耳语和诽谤，藉我们的行为而非言语成义。因为［圣经］说：\BibleQuote{多言的人，也必多听回答。善言的人就自以为义吗？从妇人生的、活的时间短的有福了：不要多言。}让我们的赞美在于天主，而不是我们自己；因为天主憎恶那些自夸的人。让他人见证我们的善行，正如我们正义的祖先那样。胆大、傲慢和狂妄属于被天主咒诅的人；而节制、谦卑和温良属于被祂祝福的人。\par

% structure: p0063 source=h2 target=section reason=relative-heading-rank; base=h2

\section{第三十一章　让我们看看藉着什么方法我们可以获得天主的祝福。}

让我们紧贴祂的祝福，并思考拥有它的方法。让我们回想从起初发生的事。我们的祖先亚巴郎为什么蒙福？岂不是因为他藉着信德行了正义和真理吗？\BibleRef{雅 2:21}依撒格怀着完全的信心，仿佛知道将要发生的事，欣然将自己献为祭品。\BibleRef{创 22:6}雅各伯因他兄弟的缘故，谦卑地离开自己的土地，来到拉班那里服事他；并且赐给了他以色列十二支派的权杖。\par

% structure: p0065 source=h2 target=section reason=relative-heading-rank; base=h2

\section{第三十二章　我们成义不是凭自己的行为，而是藉着信德。}

任何人若坦诚地细察每一件事，就会认识到从他那里所赐予的恩赐的伟大。因为从他产生了司祭和所有在祭台上服事的上主的肋未人。从他［按肉身］也产生了我们的主耶稣基督。\BibleRef{罗 9:5}从他产生了犹大支派的君王、王侯和统治者。他的其他支派也并非小有光荣，因为天主曾应许说：\BibleQuote{你的后裔将要像天上的星辰。}因此，所有这些人都受到极大的尊荣，成为伟大，不是因他们自己，也不是因他们自己的行为或所行的正义，而是因祂旨意的运作。我们也是如此，在基督耶稣内被祂的旨意所召，不是因我们自己，也不是因我们自己的智慧、理解、虔诚或发自纯洁之心的行为而成义，而是因那信德——起初全能的天主就是藉此使所有人成义；愿光荣归于祂，直到永远。阿们。\par

% structure: p0067 source=h2 target=section reason=relative-heading-rank; base=h2

\section{第三十三章　但我们不可放弃善行与爱德的实践。天主亲自为我们树立了善行的榜样。}

那么，弟兄们，我们该做什么？我们会在行善上懈怠，停止爱德的实践吗？愿天主禁止我们采取这样的行径！反而让我们以全部的精力和乐意的心情，加速完成每一项善工。因为万物的创造者和上主自己也在祂的工程中欢悦。因为藉着祂无限的大能，祂建立了诸天，并以祂不可测度的智慧妆饰了它们。祂也将大地与环绕它的水分开，并把它固定在祂自己旨意的不动根基上。祂也以祂的言辞命令地上的动物出现。同样，当祂形成海洋和其中的生物时，祂以自己的能力将它们围在［各自的界限］内。超越一切，祂以圣洁无瑕的双手形成了人——最卓越的［受造物］，并因所赋予的理解而真正伟大——祂自己肖像的活脱形象。因为天主这样说：\BibleQuote{让我们按我们的肖像，按我们的模样造人。天主于是造了人；造了男人和女人。}\BibleRef{创 1:26}完成这一切后，祂认可了它们，祝福了它们，并说：\BibleQuote{你们要生育繁殖。}\BibleRef{创 1:28}我们看到，所有的义人是如何以善行装饰，而上主自己如何以祂的工程装饰自己并欢悦。因此，既然有这样的榜样，让我们毫不迟延地服从祂的旨意，并竭尽全力从事正义的工作。\par

% structure: p0069 source=h2 target=section reason=relative-heading-rank; base=h2

\section{第三十四章　在天主那里，善行的赏报是伟大的。让我们和谐一致地团结一起，向祂祈求那赏报。}

良善的仆人满怀信心地领取他劳作的食粮；懒惰懈怠的人却不敢直视雇主的面。因此，我们必须勤于行善；因为万物都出于祂。祂这样警诫我们说：\Quote{看，主[来了]，祂的赏报在祂面前，要照各人的工作还报各人。}祂因此劝勉我们全心全意关注此事，使我们在一切善工上不致懒惰懈怠。愿我们的夸耀和信赖都在于祂。让我们顺服祂的旨意。让我们想想祂的众天使，看他们如何时刻准备着承行祂的旨意。因为圣经说：\Quote{有万万万侍立在祂周围，有千千万万服事祂，\BibleRef{达 7:10} 并呼喊说：圣、圣、圣，\OriginalTerm{万军的上主}{Lord of Sabaoth}，整个大地充满祂的光荣。}\BibleRef{依 6:3} 因此，让我们同心合意地聚集，诚恳地以同一口舌向祂呼求，好使我们能分享祂伟大而光荣的恩许。因为[圣经]说：\Quote{天主为爱祂的人所预备的，是眼所未见，耳所未闻，人心所未想到的。}\BibleRef{格前 2:9}\par

% structure: p0071 source=h2 target=section reason=relative-heading-rank; base=h2

\section{第三十五章 这赏报何其浩大。我们如何获得它？}

亲爱的弟兄们，天主的恩赐是何等有福而奇妙！不朽生命中的永生，正义中的荣光，完全信赖中的真理，确证中的信德，圣洁中的节制！这一切[如今]都为我们理智所领悟；那么，为等候祂的人所预备的，将是怎样的事呢？万世的创造者与圣父，\OriginalTerm{至圣者}{Most Holy}，唯有祂知道其数量与华美。因此，让我们竭力争取被列入等候祂的人的行列，好使我们能分享祂所应许的恩赐。但是，亲爱的弟兄们，这事要如何成就呢？如果我们的理智因信德而专注于天主；如果我们诚恳地寻求祂所喜悦和悦纳的事；如果我们做合乎祂无玷旨意的事；如果我们跟随真理之路，抛弃一切不义和邪恶，连同一切贪婪、纷争、恶行、诡诈、诽谤、毁谤，一切仇恨天主、骄傲自大、虚荣和野心。因为行这些事的人是天主所憎恶的；不仅行这些事的人，连那些赞同行这些事的人也是如此。\BibleRef{罗 1:32} 因为圣经说：\Quote{但天主对恶人说：你怎敢传述我的诫命，将我的盟约挂在嘴上？你既然憎恶管教，把我的话抛在脑后。你见了盗贼，就与他同伙，又与奸淫者分赃。你的口出恶言，你的舌编造诡诈。你坐着毁谤你的兄弟，污蔑你母亲的儿子。你行了这些事，我默然不语；你便以为我像你一样。但我要责备你，将你的罪状摆在你眼前。忘记天主的人哪！你们要思想这事，免得我像狮子撕裂你们，无人搭救。\OriginalTerm{赞颂的祭献}{the sacrifice of praise}将光荣我，有一条路，我要由此向他显示天主的救恩。}\par

% structure: p0073 source=h2 target=section reason=relative-heading-rank; base=h2

\section{第三十六章 一切福分都藉着基督赐予我们。}

亲爱的弟兄们，这就是我们找到救主的路，即是耶稣基督，我们一切祭献的\OriginalTerm{大司祭}{High Priest}，我们软弱中的护卫和帮助者。藉着祂，我们仰望高天。藉着祂，我们如同在镜中看见祂那\OriginalTerm{无玷的}{immaculate}至圣面容。藉着祂，我们心中的眼目被开启。藉着祂，我们愚昧昏昧的理智焕然一新，趋向祂奇妙的光明。藉着祂，主愿意我们尝受不朽的知识，\Quote{祂是天主荣耀的光辉，因祂所承受的名比天使更尊贵，远超过天使。}\BibleRef{希 1:3} 因为经上这样记载：\Quote{祂使祂的使者为风，使祂的仆役为火焰。}但论到祂的儿子，主这样说：\Quote{你是我的儿子，我今日生了你。你向我求，我就将外邦人赐你为基业，将地极赐你为田产。}祂又对祂说：\Quote{你坐在我的右边，等我使你的仇敌作你的脚凳。}但谁是祂的仇敌呢？就是一切恶人和那些反对天主旨意的人。\par

% structure: p0075 source=h2 target=section reason=relative-heading-rank; base=h2

\section{第三十七章 基督是我们的元帅，我们是祂的士兵。}

因此，弟兄们，让我们以全部精力，按祂神圣的诫命，履行士兵的职责。让我们想想那些在将帅麾下服役的人，他们以何等秩序、服从和顺从执行所命令的事。并非人人都是长官，或千夫长、百夫长、五十夫长或类似职位，而是各在自己的等级中执行君王和将帅所命令的事。大的不能没有小的，小的也不能没有大的。万物中都有一种混合，由此产生互助的益处。让我们以我们的身体为例。头没有脚就什么也不是，脚没有头也什么也不是；是的，我们身体上最微小的肢体对全身也是必要和有用的。但所有肢体都和谐地合作，在一条共同的法则下为保全整个身体而工作。\par

% structure: p0077 source=h2 target=section reason=relative-heading-rank; base=h2

\section{第三十八章 教会的成员当彼此顺服，谁也不可自高超过别人。}

因此，愿我们整个身体在耶稣基督内得以保全；并愿每人按所赐予他的特殊恩赐，服从邻人。强壮者不可轻视软弱者，软弱者当尊敬强壮者。富者应为穷人提供所需；穷人当赞颂天主，因为祂赐予他一个能满足其需要的人。智者当以善行而非徒有言语显其智慧。谦卑者不可自证，而应让别人为他作证。见：\BibleRef{箴 27:2} 凡在肉身内纯洁的人，不可因此自傲夸耀，应知是另有一人赐予了他节制的恩赐。那么，弟兄们，让我们思想我们是由什么材料造成的——我们来到世界时是什么样子，如同出自坟墓和完全的黑暗。那创造并塑造我们的祂，在我们出生之前就已为我们预备了丰厚的恩赐，然后将我们引入祂的世界。因此，既然我们从祂领受这一切，我们就应为一切事感谢祂；愿光荣归于祂，至于无穷世。阿们。\par

% structure: p0079 source=h2 target=section reason=relative-heading-rank; base=h2

\section{第三十九章　没有理由自满}

愚昧无思想的人，既无智慧又无教诲，嘲笑讥讽我们，一心要在自己的妄想中自高自大。因为一个必死的人能做什么？由尘土所造的人有何力量？正如经上所记载：\Quote{我眼前没有形状，只听见声音说：那么，人能在主面前洁净吗？这样的人能算为行为的纯洁吗？因为祂不信任自己的仆人，并指责祂的天使为乖谬。诸天在祂眼中也不洁净，何况那些住在土屋里的，我们自己也由之而造的呢！祂击打他们如飞蛾；从早到晚他们不能存留。因为他们不能自助，就灭亡了。祂向他们吹气，他们就死了，因为他们没有智慧。但你可以呼叫，看是否有人回答你，或者你转向任何一位圣天使；因为愤怒毁灭愚昧人，嫉妒杀死迷途者。我见过愚昧人扎了根，但他们的居所随即被吞灭。愿他们的儿子远离救恩；愿他们在比自己小的人门前被藐视，且无人拯救。因为那为他们预备的，义人要吃；他们不能脱离灾祸。}\par

% structure: p0081 source=h2 target=section reason=relative-heading-rank; base=h2

\section{第四十章　让我们在教会中保持天主所定的秩序}

既然这些事对我们显明，并且我们探究了神圣知识的深处，我们就当按主所命令我们在规定时间实行的一切，按正当秩序行事。祂规定了应献的祭品和应行的敬礼，且不是随意或杂乱地，而是在指定的时间和时辰。何处并由何人做这些事，祂自己以祂的最高旨意已予确定，为叫一切都按祂的圣意虔诚地举行，得以蒙祂悦纳。因此，凡在指定时间献上祭品的人，蒙受接纳和祝福；因为他们遵守主的律法，就不犯罪。因为大司铎有他特定的职务，司铎被指定了合适的位置，肋未人有其特别的事奉。平信徒受限于平信徒的律法。\par

% structure: p0083 source=h2 target=section reason=relative-heading-rank; base=h2

\section{第四十一章　同一主题的延续}

弟兄们，你们各人要按自己的秩序向天主感恩，以完全的良心、相称的庄重生活，不超越所规定给你们的职务规矩。弟兄们，并非在所有地方都献日常的祭、和平祭、赎罪祭和赎愆祭，而只在耶路撒冷。即使在那里，也不是在任何地方献，而只在圣殿前的祭台上，且所献之物先由大司铎和上述的侍奉者仔细查验。因此，凡做任何超出祂旨意之事的人，要受死刑。弟兄们，你们看，赐予我们的知识越大，我们所面临危险也越大。\par

% structure: p0085 source=h2 target=section reason=relative-heading-rank; base=h2

\section{第四十二章　教会中圣职人员的秩序}

宗徒们由主耶稣基督向我们宣讲了福音；耶稣基督是由天主而来。基督由天主所派遣，宗徒由基督所派遣。这两项任命都是按照天主的旨意，有秩序地完成的。他们领受了命令，并因我们主耶稣基督的复活而完全确信，又在天主的话上得以坚固，满怀圣神的保证，便出去宣告天主的国临近了。他们这样走遍城乡宣讲，便首先藉着圣神考验过的人，立为那些以后相信之人的主教和执事。这并非新事，因为早在许多世代以前，经上就已论及主教和执事。因为经上在某处这样说：\Quote{我要以正义立他们的主教，以信德立他们的执事。}\par

% structure: p0087 source=h2 target=section reason=relative-heading-rank; base=h2

\section{第四十三章　古代梅瑟平息了关于司铎尊位之争}

若那些在基督内、由天主受此托付之人，任命了上述职员，这有什么可惊奇的呢？当有福的\Person{梅瑟}，\Quote{在他全家中为忠信的仆人}，也将所受的一切命令记于圣书中，而其他先知也跟随他，一致为他的规定作证时，又当如何？因为当司铎职的争执兴起，支派间彼此争论谁堪当此荣衔时，他命令十二支派的首领将各自的杖带来，每根刻上支派的名字。他把杖取来，捆在一起，用支派首领的戒指封住，放在见证会幕内天主的桌子上。他关上会幕的门，又封了钥匙，如同封杖一样，对他们说：弟兄们，那根杖要发芽的支派，就是天主所拣选，以执行司铎职并为祂服务的。到了早晨，他召集了全以色列六十万人，向支派首领展示封印，打开见证会幕，取出杖来。\Person{亚郎}的杖不仅发了芽，而且结出了果实。亲爱的，你们怎么看？\Person{梅瑟}预先不知道这事会发生吗？毫无疑问他知道。但他这样做，是为了以色列中不产生骚乱，并使独一真天主的名受光荣。愿光荣归于祂，至于无穷之世。阿们。\par

% structure: p0089 source=h2 target=section reason=relative-heading-rank; base=h2

\section{第四十四章. 宗徒的规定：为使司铎职不再有争端。}

我们的宗徒们，藉我们的主\Person{耶稣基督}，也知道关于主教职会有纷争。因此，既然他们对此有完全的预知，就任命了上述职员，后来又指示，当这些人安息后，其他经考验的人应继承他们的职务。所以我们认为，凡由他们或后来由其他杰出人士、经全教会同意而任命的人，并谦卑、和平、无私地服务基督的羊群，且长久获得众人好评的，不应不义地免去他们的职务。若我们将那些无瑕而圣洁地履行主教职责的人驱逐，我们的罪就不算小。那些长老是何等有福，他们在此前已走完路程，获得了丰硕而完美的离去；他们不必担心有人会剥夺他们现在所被指定的位置。但我们看到，你们将一些品行优良的人从职务上免去，而这职务他们曾无瑕并光荣地履行。\par

% structure: p0091 source=h2 target=section reason=relative-heading-rank; base=h2

\section{第四十五章. 恶人迫害义人，是其本分。}

弟兄们，你们好争辩，并对与救恩无关的事充满热忱。请仔细查考圣经，那是圣神的真实言语。观察其中所写的，没有不义或虚假的内容。你们在其中找不到义人被圣洁之人所弃绝；义人确实受迫害，但那是由恶人所致；他们被投入监牢，但那是被不圣洁之人；他们被石头砸死，但那是被悖逆之人；他们被杀害，但那是被可咒骂之人，以及那些对他们怀有不义嫉妒之心的人。他们受了这些苦，却荣耀地忍受了。我们还能说什么呢？弟兄们，\Person{达尼尔}被丢入狮子坑，难道是敬畏天主的人所作的吗？\BibleRef{达 6:16} \Person{阿纳尼雅}、\Person{阿匝黎雅}和\Person{米沙耳}被投入火窑，难道是那些事奉至高者之伟大光荣敬礼的人所作的吗？\BibleRef{达 3:20} 这样的事离我们远些！那么，行这些事的是谁呢？是可憎恶的、充满一切邪恶的人，他们被激怒到如此程度，以致折磨那些以圣洁无瑕的心意事奉天主的人，不知道至高者是所有以纯洁良心敬拜祂至美之名的人的保护者和护佑者；愿光荣归于祂，至于无穷之世。阿们。然而，那些自信忍受了这些的人，如今成了光荣与尊贵的嗣子，并被天主高举、显扬，他们的纪念永世长存。阿们。\par

% structure: p0093 source=h2 target=section reason=relative-heading-rank; base=h2

\section{第四十六章. 让我们依附义人：你们的纷争是有害的。}

因此，弟兄们，我们理当追随这样的榜样；因为经上记载：\Quote{亲近圣善者，因为亲近他们的人也将成为圣善的。}又在另一处，经上说：\Quote{与无害的人相处，你必显得无害；与特选的人相处，你必成为特选的；与乖张的人相处，你必显得乖张。}因此，让我们亲近无辜和正义的人，因为他们是天主所特选的。为什么你们当中有纷争、骚动、分裂、派系和战争？我们岂不是同一的天主和同一的基督？岂不是同一的恩宠的圣神倾注在我们身上？岂不同有一在基督内的召叫？\BibleRef{弗 4:4} 我们为什么分裂并撕碎基督的肢体，对自己的身体兴起纷争，乃至疯狂到忘记\BibleQuote{我们彼此互为肢体}？\BibleRef{罗 12:5} 请记住我们的主\Person{耶稣基督}的话，祂曾说过：\Quote{祸哉，那（因他而绊倒）的人！他倒不如从未出生，免得在特选者前放置绊脚石。是的，他倒不如脖子上挂着磨石，沉入海底，也不该在小子中一人前放置绊脚石。}你们的裂教已经颠覆了许多人的信德，使许多人沮丧，使许多人怀疑，并给我们所有人带来悲伤。你们的叛乱仍在继续。\par

% structure: p0095 source=h2 target=section reason=relative-heading-rank; base=h2

\section{第四十七章. 你们近来的纷争比\Person{保禄}时代的那次更严重。}

请拿起蒙福的宗徒保禄的书信。当福音初次开始传讲时，他给你们写了什么？实在，他受圣神感动，给你们写了关于他自己、刻法和阿颇罗的事，因为那时在你们中间已经有了党派。但是那种偏爱一人的倾向使你们罪过较轻，因为你们当时的偏向是对着已经享有盛誉的宗徒，以及他们所认可的人。但现在请反省，是谁使你们偏离正道，并减少了你们闻名遐迩的兄弟之爱的声誉？亲爱的弟兄们，这真是可耻，极其可耻，且与你们的基督徒身份不相称：竟会听到这样的事，即最坚定、最古老的格林多教会，因为一两个人，就对其长老们进行叛乱。这个谣言不仅传到了我们这里，也传到了那些与我们无关的人那里；因此，由于你们的愚妄，主的名被亵渎，同时危险也临到你们自己身上。\par

% structure: p0097 source=h2 target=section reason=relative-heading-rank; base=h2

\section{第四十八章　让我们回到实践兄弟之爱。}

因此，让我们赶紧结束这种状况；让我们跪在主前，含泪恳求祂，求祂仁慈地与我们和好，并恢复我们从前合宜而圣洁的兄弟之爱。因为这种行为是正义之门，为了获得生命而敞开的，正如经上记载：\Quote{给我敞开正义之门，我要进去赞美主；这是主的门，义人要从它进去。}因此，虽然许多门已经敞开，但这正义之门是在基督内的门，所有进入并行走在圣洁和正义中、一切按秩序行事的人都是有福的。一个人可以有信德；可以善于表达知识；可以有智慧判断言语；可以在一切行为上纯洁；然而，他越在这些方面显得比别人优越，就越应当谦逊，并寻求众人的共同益处，而不是仅仅自己的利益。\par

% structure: p0099 source=h2 target=section reason=relative-heading-rank; base=h2

\section{第四十九章　爱的赞颂。}

让那在基督内有爱的人遵守基督的诫命。谁能描述天主之爱的有福联系？谁能恰当地述说其美的卓越？爱所提升的高度是不可言喻的。爱使我们与天主结合。爱遮盖许多罪恶。爱包容一切，在一切事上忍耐。爱中没有卑鄙，没有傲慢。爱不容分裂；爱不引起叛乱；爱在和谐中行一切事。天主的全体选民都因爱而成全；没有爱，什么也不能使天主喜悦。主因爱接纳了我们。由于祂对我们的爱，我们的主耶稣基督按照天主的旨意为舍了自己的血；祂的肉身为了我们的肉身，祂的灵魂为了我们的灵魂。\par

% structure: p0101 source=h2 target=section reason=relative-heading-rank; base=h2

\section{第五十章　让我们祈求被认为配得爱。}

亲爱的弟兄们，你们看爱是多么伟大而奇妙的事，它的完美简直无法言说。谁配得上它呢？除非天主使某人如此。因此，让我们祈求并恳求祂的仁慈，使我们能在爱中无过犯地生活，摆脱一切对人的偏爱。从亚当直到今日的所有世代都已过去；但是那些藉着天主的恩宠在爱中得以成全的人，现在已在义人之中拥有一个位置，并将在基督之国显现时显明出来。因为经上记载：\BibleQuote{进入你的内室，稍待片刻，直到我的怒火过去；然后我会记起施恩的日子，并将你从坟墓中举起。}\BibleRef{依 26:20}亲爱的弟兄们，如果我们和谐地遵守天主的诫命，我们就有福了；这样，藉着爱，我们的罪得以赦免。因为经上记载：\Quote{过犯得赦免、罪恶被掩盖的人是有福的。主不将罪归于他、口中没有诡诈的人是有福的。}这福分临到那些藉着我们的主耶稣基督被天主拣选的人；愿光荣归于祂，直到永远。阿们。\par

% structure: p0103 source=h2 target=section reason=relative-heading-rank; base=h2

\section{第五十一章　让参与纷争的人承认他们的罪。}

因此，让我们为那些因仇敌的诱惑而犯下的一切过犯恳求宽恕。这些曾是叛乱和分歧的领袖的人，应当关注共同的希望。因为那些生活在敬畏和爱中的人，宁愿自己受苦，也不愿邻人受苦。他们宁愿自己承担责备，也不愿我们良好而虔诚传承下来的和睦受到损害。因为人承认自己的过犯，比硬着心肠更好，就像那些在梅瑟天主的仆人面前煽动叛乱的人的心肠那样刚硬，他们的定罪已显明给众人。因为他们活着下了阴府，死亡吞没了他们。法郎和他的军队，以及埃及的众首领和战车骑兵，都沉入红海深处而灭亡，没有别的原因，只因他们愚昧的心刚硬，尽管在天主的仆人梅瑟手中，在埃及地行了那么多神迹奇事。\EditorialNote{出谷纪第十四章}\par

% structure: p0105 source=h2 target=section reason=relative-heading-rank; base=h2

\section{第五十二章　这样的宣认中悦天主。}

弟兄们，主一无所缺；除了向他宣认之外，他不向任何人要求什么。因为蒙选的达味说：\Quote{我要向主宣认，这比献上带角有蹄的公牛更令祂喜悦。愿贫苦的人看见并欢喜。}他又说：\Quote{你要向上主献上赞美的祭，向至高者还你的愿。在患难之日呼求我，我必拯救你，你也要光荣我。}因为\Quote{天主所要的祭是痛悔的精神。}\par

% structure: p0107 source=h2 target=section reason=relative-heading-rank; base=h2

\section{第五十三章　梅瑟对其子民的爱。}

亲爱的，你们明白，你们确实明白圣经，并且非常认真地查考了天主的圣言。所以要将这些事记在心里。当梅瑟上山，在那里禁食谦卑，四十昼夜，主对他说：\Quote{梅瑟，梅瑟，快快从这里下去，因为你从埃及地领出来的百姓已经犯了罪。他们很快偏离了我命令他们走的道路，为自己铸造了偶像。}主又对他说：\Quote{我曾一再对你说：我看见了这百姓，看哪，他们是硬着颈项的百姓。让我消灭他们，从天下涂抹他们的名；我要使你成为一个强大而奇妙的民族，比这更众多。}但梅瑟说：\BibleQuote{主，万万不可！求你赦免这百姓的罪，否则也把我从生命册上涂抹吧。}\BibleRef{出 32:32} 何等奇妙的爱！何等无与伦比的成全！仆人对他的主说话无拘束，为百姓求赦免，或祈求自己也与他们一同灭亡。\par

% structure: p0109 source=h2 target=section reason=relative-heading-rank; base=h2

\section{第五十四章 凡充满爱德者，为使教会重获平安，愿承受一切损失。}

你们中间谁有高贵的心？谁有同情心？谁充满爱德？让他宣告：\Quote{如果因我而起了叛乱、不和与分裂，我愿意离开，我愿意到你们愿意的任何地方去，并照多数人的吩咐去做；只愿基督的羊群与治理他们的司铎和平相处。}如此而行的人，必在主内为自己赢得极大的荣耀；每个地方都将欢迎他。因为\Quote{大地和其中的万物属于主。}这些事，凡度着永不后悔的虔敬生活的人，已经做了，也常会做。\par

% structure: p0111 source=h2 target=section reason=relative-heading-rank; base=h2

\section{第五十五章 此类爱德的典范。}

为引进一些外邦人的例子：许多君王和王子，在瘟疫时期，得到神谕的指示，献出自己的生命，好以自己的血拯救他们的同胞脱离毁灭。许多人离开自己的城邑，为平息城内的叛乱。我们中间也有许多人为赎回他人而自陷缧绁。还有许多人卖身为奴，用自己所得的钱来供养他人。还有许多妇女，因天主的恩宠而坚强，做出了许多勇士般的壮举。蒙福的友弟德，当她的城被围困时，她向长老们请求许可进入外人的营寨；冒险出去，为了她对被围的祖国和人民的爱；主将敖罗斐乃交在一个妇人手中。\BibleRef{友 8:30} 艾斯德尔也是信德完备，为了拯救以色列十二支派脱离迫近的毁灭，冒了同样的危险。因为她以禁食和谦卑恳求那鉴察万物的永恒天主；天主看到她心灵的谦卑，便拯救了那因她而面临危险的百姓。\par

% structure: p0113 source=h2 target=section reason=relative-heading-rank; base=h2

\section{第五十六章 让我们彼此劝诫和纠正。}

那么，我们也要为那些陷于任何罪中的人祈祷，求赐给他们温良和谦卑，好使他们服从，不是服从我们，而是服从天主的旨意。因为这样，他们将从我们这里获得丰硕而完美的纪念，我们对他们的同情，既在我们向天主的祈祷中，也在我们向圣人们的提及中。亲爱的，让我们接受纠正，谁也不应因此不悦。我们彼此劝诫的那些教训，本身是好的，也极有裨益，因为它们使我们与天主的旨意合一。因为圣言这样说：\Quote{主严厉地管教我，却没有将我交于死亡。}\Quote{主所爱的，祂必管教，又鞭打凡所收纳的儿子。}\Quote{义人，}它说，\Quote{要在仁慈中管教我，并责备我；但不要让罪人的油抹在我的头上。}他又说：\BibleQuote{蒙主责备的人是有福的，不可拒绝全能者的训诲。因为祂使人忧伤，又使人重获喜乐；祂击伤，祂的手又医治。在六次患难中，祂必拯救你；在第七次，灾祸也不会触及你。饥荒中，祂要救你脱离死亡；战争中，祂要释放你脱离刀剑的权势。祂必隐藏你免受舌头的鞭打，灾祸来临时你也不惧怕。你要嗤笑不义和邪恶的人，田野的野兽你也不用害怕。因为野兽必与你和平相处：那时你便知道你的家宅平安，你的帐幕住所不会倒塌。你也要知道你的后裔众多，你的子孙如田野的草。你必如按时收割的熟谷，或按食聚集的禾捆，进入坟墓。}\BibleRef{约 5:17} 亲爱的，你们看，\Quote{受主管教的人必得保护；因为天主是良善的，祂纠正我们，好使我们因祂神圣的管教而受到劝诫。}\par

% structure: p0115 source=h2 target=section reason=relative-heading-rank; base=h2

\section{第五十七章 愿制造纷争者服从。}

所以你们这些惹起这纷争的，当服从司铎，接受纠正，好能悔改，屈膝于心。学习服从，抛弃你们舌头那骄傲和狂妄的自恃。因为你们在基督的羊群中占据卑微却荣耀的位置，胜过自高自大而被逐出祂百姓的望德。因为全德的智慧这样说：\Quote{看，我要向你们说出我圣神的话语，并教导你们我的言语。因为我呼唤，你们却不听；我伸出我的话语，你们却不理会，反而藐视我的劝告，不听从我的责备；因此我也要嘲笑你们的毁灭；是的，当灾难临到你们，当突然的混乱笼罩你们，当倾覆如旋风般显现，或苦难与压迫加于你们时，我必喜乐。因为将来你们呼求我，我不听；恶人寻找我，却寻不见。因为他们憎恨智慧，不选择敬畏上主；也不听从我的劝告，反倒藐视我的责备。所以他们要自食其果，充满自己的不虔敬。\BibleRef{箴 1:22}……因为，作为他们向婴孩所行之恶的惩罚，他们必被杀，而查究对不虔敬者就是死亡；但听从我的，必安居于希望，不受任何邪恶的惊扰。}\par

% structure: p0117 source=h2 target=section reason=relative-heading-rank; base=h2

\section{第五十八章　服从是救恩的先声}

因此，让我们逃离智慧对悖逆者所发出的警告威胁，臣服于祂至圣至荣的圣名，好使我们能信赖祂威严的至圣尊名。接受我们的劝告，你们将无后悔。因为正如天主活着，主耶稣基督和圣神也活着——正是蒙选者的信德与望德——谁若心存谦卑，常怀温柔，毫不反悔地遵守天主所赐的命令和规定，他必在那些藉耶稣基督而得救的人中得着位置和名分，愿荣耀藉着祂归于天主，直到永远。阿们。\par

% structure: p0119 source=h2 target=section reason=relative-heading-rank; base=h2

\section{第五十九章　警告悖逆与祈祷}

但如果有人不听从祂藉我们所说的话，让他们知道他们必使自己陷入过犯和严重的危险；然而我们在这罪上无咎，我们要不断祈祷和恳求，愿万有的创造者藉祂的爱子耶稣基督，保全祂在全世界蒙选者的完整数目，祂曾藉祂从黑暗召我们到光明，从无知到认识祂荣耀圣名的知识，我们的希望寄托于祢的圣名，它就是一切受造物的首要原因——祢开启了我们的心眼，使能认识祢，唯独祢\Quote{居于至高者之上，圣者之中为圣}，\BibleRef{依 57:15}祢\Quote{贬抑骄傲者的狂傲}，\BibleRef{依 13:11}祢\Quote{摧毁异邦人的计谋}，祢\Quote{使卑微者升高，使高傲者降低}；祢\Quote{使人富足也使人贫穷}，\BibleRef{撒上 2:7}祢\Quote{使人死也使人活}，\BibleRef{申 32:39}惟独祢是众灵的恩主、一切血肉的天主，祢察看深渊，是世人行为的见证，危难者的帮助，绝望者的救主，一切灵的创造者和守护者，祢使地上万国增多，并从中拣选了藉着祢的爱子耶稣基督爱祢的人，借着他祢教导、圣化、尊荣我们。主，我们愿祢作我们的帮助和援助。求祢拯救我们中患难的人，怜悯卑微者，扶起跌倒者，向需要者显现，医治病人，引领祢百姓中迷途者，饱饫饥饿者，赎回我们中被捆绑者，扶持软弱者，安慰灰心者；愿万民都知道惟有祢是天主，耶稣基督是祢的圣子，我们是祢的子民、祢草场的羊群。\par

% structure: p0121 source=h2 target=section reason=relative-heading-rank; base=h2

\section{第六十章　祈祷续}

祢亲手的工作使世界持久的结构显现；主，祢创造了我们所居住的大地——祢世世代代信实，审判公正，能力与威严奇妙，以智慧创造，以明智固定所造之物，在得救的人中祢是善的，在信赖祢的人中祢是信实的；哦，仁慈而怜悯的一位，求祢宽恕我们的不义、过犯、悖逆与罪愆。不要记算祢仆人和使女的每一项罪，但求祢以真理的洁净净化我们；指引我们的脚步，使我们能存心圣洁行走，行祢眼中看为美好和喜悦的事，也在我们统治者眼前如此。是的，主，求祢使祢的圣容光照我们，赐我们平安美善，使我们得蒙祢大能的手庇护，被祢高举的臂膀拯救脱离一切罪恶，并救我们脱离无故恨我们的人。求祢赐给我们和地上所有居民和谐与平安，如同祢赐给我们列祖那样，他们曾在信德和真理中呼求祢，我们同样顺服于祢全能至上的圣名。\par

% structure: p0123 source=h2 target=section reason=relative-heading-rank; base=h2

\section{第六十一章　祈祷续——为统治者和执政者；结语}

至于我们地上的统治者和执政者——主，祢曾以荣耀而不可言喻的大能赐给他们王权，为使我们知晓祢所赐给他们的荣耀和尊荣，并服从他们，丝毫不抗拒祢的旨意；主，求祢赐给他们健康、平安、和谐、稳定，使他们能毫无过犯地行使所赐的权柄。因为祢，天上的主、永恒的君王，将荣耀、尊荣和掌管地上事物的权柄赐给世人；主，求祢按祢眼中看为美好和喜悦的事指引他们的计谋，使他们能虔诚地、和平地、温柔地行使祢所赐的权柄，从而得蒙祢的悦纳。哦，惟独祢有能力成就这些事，并为我们成就更丰盛的美善，我们藉着我们灵魂的大司祭和守护者耶稣基督赞美祢，愿荣耀和威严藉着祂归于祢，从今时直到世世代代，直至永远。阿们。\par

% structure: p0125 source=h2 target=section reason=relative-heading-rank; base=h2

\section{第六十二章  总结与结论——论虔敬。}

诸位弟兄，关于那些与我们敬礼相关的、对追求虔敬与公义生活的人最有益的事，我们已经足够详尽地写信给你们。因为论及信德、悔改、真爱、节制、清醒和忍耐，我们已触及每一段经文，提醒你们应当以正义、真理和宽容，在圣洁中、一心一意、不怀怨恨、以爱德与和平、且以恒久的温良，中悦全能的天主，正如我们先前所提到的祖先，因他们思想上的谦卑，蒙悦纳于圣父、天主、造物主及全人类。我们更乐意提醒你们这些事，因为我们确信我们是在写信给忠信、德高望重、并曾深究天主训诲神谕的人。\par

% structure: p0127 source=h2 target=section reason=relative-heading-rank; base=h2

\section{第六十三章  劝勉——信件由特使送出。}

因此，效法如此众多且如此良好的榜样，是合宜的；并要低下头颈，履行服从的本分，以便我们不被虚妄的纷争扰乱，在真理中毫无瑕疵地达到摆在我们面前的目标。如果你们服从我们所写的这些话，并藉着圣神，按照我们在这封信中为和平与合一所作的恳求，根除你们嫉妒中不合法的忿怒，你们就将使我们充满喜乐与欢欣。我们已派遣了忠信而明智的人，他们在我们中间从幼年到老年的行为都是无可指摘的——这些人将在你们和我们之间作证。我们这样做，是要你们知道，我们全部的关心一直是，并且仍然是，使你们尽快归于和平。\par

% structure: p0129 source=h2 target=section reason=relative-heading-rank; base=h2

\section{第六十四章  为所有呼求天主的人祈求祝福。}

愿那洞察万有、统辖万灵、主宰一切血肉的天主——祂拣选了我们的主耶稣基督，并藉着祂拣选了我们成为特殊的子民\BibleRef{铎 2:14}——赐给每一个呼求祂荣耀而圣洁之名的灵魂：信德、敬畏、和平、忍耐、宽容、节制、纯洁和清醒，以中悦祂的名，藉着我们的至高司祭和保护者耶稣基督；愿荣耀、威严、权能、尊崇，藉着祂归于祂，从今时直到永远。阿们。\par

% structure: p0131 source=h2 target=section reason=relative-heading-rank; base=h2

\section{第六十五章  劝勉格林多人速速回报和平已恢复。祝福。}

请速速平安喜乐地给我们派回这些到你们那里的使者：\Person{克劳狄·埃斐布斯}、\Person{瓦莱里·比托}以及\Person{福图纳图斯}；好让他们能更早地向我们宣布我们所热切渴望和期盼的［在你们中间的］和平与和谐，使我们能更快地为你们中间恢复的良好秩序而欢欣。愿我们主耶稣基督的恩宠与你们同在，也与所有藉着祂蒙天主召叫、从永远到永远的各处的人同在。愿荣耀、尊崇、权能、威严和永恒的统治，藉着祂归于祂，从永远到永远。阿们。\par

\SourceRecord{Translated by John Keith.；Ante-Nicene Fathers；Vol. 9.；Edited by Allan Menzies.；Buffalo, NY: Christian Literature Publishing Co.,；1896.；<http://www.newadvent.org/fathers/1010.htm>.}
